\newcommand{\onlinesummary}{Learning from a fitted prior adds little: on LongMemEval-S recall moves from \V{on-lme-static} to \V{on-lme-online}, on LoCoMo from \V{on-locomo-static} to \V{on-locomo-online}, and on RepoMem from \V{on-repomem-static} to \V{on-repomem-online}. Learning matters when the prior is wrong: from a flat prior, static recall is \V{on-lme-cold-static} on LongMemEval-S and \V{on-locomo-cold-static} on LoCoMo, and the online policy recovers to \V{on-lme-cold-online} and \V{on-locomo-cold-online}, close to the fitted prior. The UCB bonus changes little. Noise is the danger. With 20\% of rewards flipped, online learning does \emph{worse} than not learning at all: \V{on-lme-online-noisy} on LongMemEval-S, \V{on-locomo-online-noisy} on LoCoMo and \V{on-repomem-online-noisy} on RepoMem, because the flips raise the weight of the many irrelevant items. The standard correction for a known flip rate recovers part of the loss on LongMemEval-S and most of it elsewhere (\V{on-lme-online-debiased}, \V{on-locomo-online-debiased} and \V{on-repomem-online-debiased}), but a real deployment does not know its flip rate. So we ship policy learning as opt-in. We did not test whether a larger ridge would make it safe.}
\newcommand{\consolidationsummary}{Consolidation is not free. At $\tau=0.9$ it merges only near-duplicates: it saves \V{consave-lme-0.9-0.5}\% of tokens on the LongMemEval-S sample and \V{consave-locomo-0.9-0.5}\% on LoCoMo, with recall unchanged. At $\tau=0.8$ it saves \V{consave-lme-0.8-0.5}\% and \V{consave-locomo-0.8-0.5}\%, and recall moves from \V{con-lme-None-0.0} to \V{con-lme-0.8-0.5} and from \V{con-locomo-None-0.0} to \V{con-locomo-0.8-0.5}. At $\tau=0.7$ the loss is larger. Since the budget caps injected tokens anyway, a smaller store buys little here, so the default threshold is $0.9$.}
\newcommand{\forgettingsummary}{Forgetting has a real cost on this benchmark. Idle rules of 365 and 180 days cut the active store to \V{fgt-idle-365d} and \V{fgt-idle-180d} tokens from \V{fgt-keep-all}, as Theorem~\ref{thm:store} leads one to expect, but recall falls from \V{fg-keep-all} to \V{fg-idle-365d} and \V{fg-idle-180d}, because long-lived facts are read rarely and are still correct. Sessions in this benchmark are \V{gapdays} days apart, so these idle rules are short in session terms, and the numbers do not transfer directly to a daily workflow. Token capacities behave similarly. Given the latencies below, a store of this size does not need forgetting, so the rules are opt-in. Archived items remain recoverable by hand, but they leave retrieval.}
\newcommand{\conclusionsummary}{On retrieval metrics, cost-aware selection raises recall per token where item costs vary and relevant items are cheap. A cross-encoder remains the stronger scorer at small budgets, and the objective can absorb it as a feature. Stale injection is a large problem on code facts, and simple rules fix most of it, so the probabilistic validity model has not yet shown value beyond them. Free-text updates remain unsolved. Plain hybrid stays better at large budgets on LoCoMo, and online learning helps mainly when the prior is wrong and is fragile under reward noise.}
